\problem{ Estimation de la température de la Terre }

\question{À partir des observations précédentes, estimez la température de la surface de la Terre.}

La puissance du Soleil est:
\[
    P_{Soleil} = 4\pi R_{Soleil}^2 \sigma T_{Soleil}^4. \notag
\]

Puisque la Terre est assez loin du Soleil, l'intensité du Soleil perçu par la Terre est seulement représentée 
par des rayons qui arrivent à la Terre parallèlle entre eux. Nous avons donc:
\[
    I = \frac{P_{Soleil}}{4\pi d^2}. \notag
\]

La puissance absobée par la Terre est donc :
\[
    P_{abs} &= \frac{P_{Soleil}}{4\pi d^2} \pi R_{Terre}^2 \notag \\
    &= \frac{P_S}{4 d^2} R_{Terre}^2. \notag
\]

Cette valeur étant égale à la puissance émise de la Terre, nous avons :
\[
    P_{Terre} &= P_{abs} \notag \\
    4\pi R_{Terre}^2 \sigma T_{Terre}^4 &= \frac{4\pi R_{Soleil}^2 \sigma T_{Soleil}^4}{4 d^2} R_T^2 \notag \\
    T_{Terre} &= \qty(\frac{T_{Soleil}^4 R_{Soleil}^2}{4 d^2})^{\frac{1}{4}}. \notag
\]

Avec les valeurs données, nous trouvons :
\[
    T_{Terre} &= \qty(\frac{(5500)^4 (7\times10^8)^2}{4 (1.5\times10^11)2})^{\frac{1}{4}} \notag \\
    &\approx 266K.
\]
\clearpage
